\section{三次原始を用いない四次係数}
\begin{lemma}[有限BCHの縮小代表]\label{mix:bchfour}
$M=M_0+hS_1+h^2S_2+h^3S_3+h^4S_4\pmod{h^5}$を曲率付きMaurer--Cartan元とし、
$d=[M_0,-]$、$dv=S_1,dw=S_2+[v,S_1]/2-Q_2$、$Q_2$は純曲率とする。
次のcochainは閉で、既知の二次gaugeを施した四次項と同じ曲率類を表す：
\begin{align*}
 R_4={}&S_4+[v,S_3]+\tfrac14[v,[v,S_2+Q_2]]\\
       &+\tfrac12[w,S_2+Q_2]+\tfrac12[v,[w,S_1]].
\end{align*}
$v$のarityが三、$w$が一と三、$S_r$が零と二なら最大arityは六である。
未計算の三次gaugeはこの四次類に影響しない。
\end{lemma}
\begin{proof}
$A=e^{h\operatorname{ad}v}M$とすると$A_1=0$、
$A_2=S_2+[v,S_1]/2$、
$A_4=S_4+[v,S_3]+[v,[v,S_2]]/2+[v,[v,[v,S_1]]]/8$である。
$B=e^{h^2\operatorname{ad}w}A$では$B_1=0,B_2=Q_2$、
$B_4=A_4+[w,A_2+Q_2]/2$となる。
三次gaugeの四次寄与は$[g_3,B_1]=0$である。
Maurer--Cartan式から$dB_4+[Q_2,Q_2]/2=0$、零入力同士の括弧は零なので閉である。
$P=[v,[v,w]]/4$とし、導分則で展開すると$B_4-dP=R_4$となる。
例えば
\[
 4dP=[S_1,[v,w]]+[v,[S_1,w]]+[v,[v,S_2-Q_2]]
       +[v,[v,[v,S_1]]]/2.
\]
これとJacobi恒等式を用いればよい。
この境界に最大arity八の項が吸収され、表示の最大arityは六となる。
\end{proof}

\begin{proposition}[二係数を読むJacobian商]\label{mix:fourquot}
四次の共変純曲率の可能形を
$A\sum e_ix_i^2+B\sum e_je_kx_i+C\sum e_im_i$とすると、
六級数規格化での係数は$[h^4]V_e=B-2A$、$[h^4]T_e=C-A$である。
$K'=\mathbb Q[e_1,e_2,m_1]/(e_1^2,e_2^2,m_1^2)$、$e_3=m_2=m_3=0$へ特殊化してもこの二次元商を検出できる。
\end{proposition}
\begin{proof}
$x_i=X_i-Ah^4e_i$の変化は
$-Ah^4\sum e_iF_i=-Ah^4(\sum e_ix_i^2+2\sum e_je_kx_i+\sum e_im_i)$である。
共通拡大は$1+O(h^6)$なので四次係数へ影響しない。
特殊化後の十五Jacobian列$gF_i$、$g\in\{x,y,z,e_1,e_2\}$と二列$e_1e_2z,e_1m_1$は独立である。
実際、線形結合の純三次係数からベクトル場の$x,y,z$係数が零、
次に$e_\alpha x_i^2$係数から残る係数も零となり、二検出係数も零となる。
従って特殊化はこの商上単射である。
\end{proof}

\paragraph{実行する有限再構成.}
二次原始$w_3$を全係数上で生成し、対角指標へ射影してから六置換で平均し、初めて$K'$へ特殊化する。
実際のPBW規則は
\[
 yx=xy-h(z^2+e_1y+e_2x),\quad zx=xz+h(y^2+e_1z),\quad
 zy=yz-h(x^2+e_2z+m_1)
\]
である。この語還元と完全対称化$T$、有限幾何級数の$T^{-1}$からWeyl積を得る。
正の入力微分指数$(a,b)$を$|a|+|b|$の昇順に取り、
単項式入力$x^a,x^b$への作用から既知低微分成分を引いて$a!b!$で割れば、通常微分記号の当該係数を再生成できる。
総入力次数四以下は141組である。

係数重み四以上を消す射影を$J_{\le3}$とする。
$C_r$の係数重みは全入力微分次数に等しく、$[v,-]$は既存係数を高々一回微分するので、
$J_{\le3}R_4$には$C_3$の微分次数四以下、$C_4$の三以下だけで十分である。
内側bracketでは外側で必要となる重み四までを保持し、全nested bracketの後に初めて最終射影を行う。
Hochschild微分と有限cobar縮約は係数線形であり、曲率挿入は既存係数を微分せず$F$の偏微分を掛けるため、
係数重み四以上のイデアルを保つ。
二入力HKR成分については、特殊化前の全$K$上の閉cochainを先に考え、正則勾配列と重みで$\pi_F$の定数倍と決めてから射影する。
射影後のKoszul複体そのものの完全性は仮定しない。

\paragraph{検証の論理.}
現行検証器は$C_3$の90微分型と$C_4$の24型を生成し、同じ操作から得た$C_1,C_2$を既知の全記号と照合する。
有限BCH式の閉性を全入力記号で検査し、arity六、四の順に有限poset縮約を行う。
各段で原始のHochschild微分が対象成分に等しいことを確認する。
二入力の双ベクトルを除き、最後の一入力原始を引いて純曲率を得る。
この計算は全$C_3,C_4$の無制限な表を検証する主張ではなく、証明した射影に十分な低記号を検証するものである。

\begin{theorem}[固定中心の四次量子係数]\label{mix:fourth}
定理\ref{mix:secondcurve}の二次gaugeを上の有限群平均で固定し、係数環を保つ共変な全次数帰納へ延長すると
\[
 V_e(h)=h^4/9+O(h^{10}),\qquad T_e(h)=h^4/18+O(h^{10})
\]
となる。
\end{theorem}
\begin{proof}
現行の有限再構成は$J_{\le3}R_4$を170131微分型として生成する。
全入力記号で閉性を確認し、arity六の86332型の原始を除くと3443型、
arity四の1607型の原始を除くと46型が残る。
各段で原始のHochschild微分との完全一致を確認する。
二入力のHKR成分を除き一入力原始を引くと、純曲率は厳密に
\[
 q_4=-\frac5{144}(e_1x^2+e_2y^2)+\frac1{24}e_1e_2z+\frac1{48}e_1m_1
\]
となる。これは有限の有理数多項式恒等式であり、期待される二数値は源や原始の探索条件へ入れていない。
全十七Jacobian列との一意な分解は
\[
 q_4=-\frac5{144}(e_1F_x+e_2F_y)+\frac19e_1e_2z+\frac1{18}e_1m_1.
\]
実際$e_1F_x+e_2F_y=e_1x^2+e_2y^2+2e_1e_2z+e_1m_1$である。
命題\ref{mix:fourquot}から後の二係数が全係数上の二次元商を検出する。
三次の純曲率は対角指標の許す各置換表現に符号表現がないため、群平均した有限原始で零にできる。
この三次gaugeが四次類を変えないことは補題\ref{mix:bchfour}で示した。
以後のgaugeは四次以下を保存し、六級数の合同条件から次の許容次数は十である。
\end{proof}

\paragraph{再現.}
\texttt{computations/verify\_mixed\_fourth\_order\_gauge.py}は、
二次原始の生成器を呼び、PBW/Weyl演算、有限群平均、BCH、poset縮約、
各原始恒等式と二係数の読取りを全て再生成する。
SymPyとpython-flintを用い、外部JSON、数値補間、有限体の零判定は用いない。
\texttt{verify\_mixed\_quadratic\_dual.py}は別に階数32模型の全32768結合律を$\mathbb Q[\alpha,h]$上で確認する。
